\section{Local continuation and the feedforward sublayer}\label{sec:local-continuation}

Attention refines the receiver interface $Q_j$. The pointwise feedforward sublayer occupies the other side of the central factorization: it refines the receiver-local continuation $G_j$. This distinction explains both the exact hidden-coordinate computation and its architectural limit. The FFN can construct new local observables from the retained receiver state, but it cannot recover source distinctions already identified by $Q_j$.

\subsection{Marked refinements of a local continuation}

\begin{definition}[Marked local-continuation refinement]\label{def:local-continuation-refinement}
Let $G:Z\to W$ be a receiver-local continuation. A marked refinement of $G$ through a state space $A$ consists of maps
\[
Z\xrightarrow{a}A\xrightarrow{o}W
\]
with $G=o\circ a$, together with declared projections or coordinates on $A$ that are part of the architectural presentation.
\end{definition}

The mark is necessary. Every function can be padded with arbitrary identity maps, invertible recodings, or unused coordinates. A hidden state is architectural only when the specified continuation actually contains that intermediate presentation and its coordinate roles.

\subsection{Two-layer local continuation}

Consider the receiver-local map
\begin{equation}\label{eq:ffn}
\FFN(x)=W_2\phi(W_1x+b_1)+b_2,
\end{equation}
where $\phi$ acts coordinatewise. Let
\[
a(x):=\phi(W_1x+b_1)\in\R^{d_{\mathrm{ff}}}
\]
and let $E=\{1,\ldots,d_{\mathrm{ff}}\}$ index the marked activation coordinates. Write $w_e$ for the $e$th column of $W_2$.

\begin{proposition}[Marked hidden-coordinate decomposition]\label{prop:ffn-hidden-carrier}
The specified two-layer continuation factors through the marked activation presentation
\[
x\xrightarrow{a}\R^{d_{\mathrm{ff}}}\xrightarrow{y\mapsto W_2y+b_2}W,
\]
and
\begin{equation}\label{eq:ffn-direction-sum}
\boxed{\FFN(x)=b_2+\sum_{e\in E}a_e(x)w_e.}
\end{equation}
\end{proposition}

\begin{proof}
The coordinatewise activation produces the marked vector $a(x)=(a_e(x))_e$. Multiplication by $W_2$ forms the linear combination of its columns with those coefficients.
\end{proof}

The conclusion is stronger than a formal column expansion and weaker than a mechanism claim. The hidden coordinates are architectural because the declared two-layer sublayer computes and retains them as the intermediate input to $W_2$. The proposition does not establish that one hidden coordinate is an independent expert, causal mechanism, or uniquely individuated basis direction.

Unlike attention coefficients, the native $a_e(x)$ may be signed and need not have unit sum. A nonnegative normalized allocation can be manufactured by signed carrier doubling and amplitude extraction, but that is a derived representation rather than the native form. The exact construction is recorded in the supplement.

\subsection{Gated local continuation}

Many contemporary blocks use
\begin{equation}\label{eq:gated-ffn}
\GFFN(x)
=
W_o\left[\phi(W_gx+b_g)\odot(W_ux+b_u)\right]+b_o.
\end{equation}

\begin{proposition}[Gated hidden-coordinate decomposition]\label{prop:gated-ffn}
Let $w_e$ be the $e$th column of $W_o$ and define
\[
a_e(x)
:=
\phi((W_gx+b_g)_e)(W_ux+b_u)_e.
\]
Then
\[
\boxed{\GFFN(x)=b_o+\sum_e a_e(x)w_e.}
\]
The marked intermediate carrier is the coordinatewise gated activation vector.
\end{proposition}

\begin{proof}
The vector inside $W_o$ has coordinates $a_e(x)$, and $W_o$ forms the corresponding column combination.
\end{proof}

The gated form changes the local coefficient constructor while preserving the architectural pattern: receiver-local state constructs coefficients on a marked hidden coordinate family, then a shared output map recombines the resulting directions.

\subsection{Local observable construction and its limit}

Let the receiver interface after routed exposure be
\[
Q_i(H)\cong\bigl(h_i,z_i(H)\bigr),
\]
and let a residual and normalization map produce the FFN input $x_i$. Every pointwise FFN output is then a function of $x_i$. The sublayer can:
\begin{itemize}
\item construct nonlinear observables already determined by the retained receiver macrostate;
\item move those observables into coordinates that later restricted operations can use;
\item apply one shared local law across receiver occurrences, possibly conditioned on explicit type marks.
\end{itemize}
It cannot distinguish predecessor states that induce the same $x_i$. By \Cref{thm:squared-decomposition}, increasing local approximation capacity addresses restricted accessibility but not unresolved variation inside the interface fibers.

This yields the division of labor
\[
\boxed{
\begin{aligned}
\text{source-resolved routing}
&\text{ determines which nonlocal distinctions are constructed},\\
\text{aggregation and residual retention}
&\text{ determine which receiver state survives},\\
\text{the FFN}
&\text{ constructs local observables of that retained state}.
\end{aligned}}
\]

\subsection{Coordinate nonuniqueness is activation-relative}

For a linear factorization $W=AB$, every invertible intermediate change of basis $M$ gives $W=(AM^{-1})(MB)$. A nonlinear two-layer FFN does not admit arbitrary $GL(d_{\mathrm{ff}})$ changes of hidden basis, because a general $M$ does not commute with coordinatewise $\phi$.

The valid presentation symmetries are activation-compatible. They include hidden-unit permutations, with corresponding permutations of the rows of $W_1$ and columns of $W_2$. For positively homogeneous activations, suitable positive coordinate rescalings are also available when compensated between the incoming and outgoing weights. Further symmetries require separate proof for the chosen activation and parameterization. The paper therefore does not identify all algebraic factorizations of the external local map with the specified hidden architecture.

\subsection{A derived hidden potential is not foundational}

If one first doubles signed coordinates and normalizes their nonnegative magnitudes, the resulting distribution can be written in Gibbs form on its positive support by defining $-\tau_{\mathrm{hid}}\log p_x$. This is tautologically exact for any positive finite distribution. It is weaker than the attention result because the native FFN coefficients need not arise from a positive compositional evidence law, preserve one receiver-relative ratio structure, or place total amplitude inside the normalized row. The useful commonality is state-conditioned coefficient construction over a marked carrier, not a shared thermodynamic origin.
